\documentclass[conference]{IEEEtran}
\usepackage{cite}
\usepackage{amsmath,amssymb,amsfonts}
\usepackage{graphicx}
\usepackage{booktabs}
\usepackage{url}
\usepackage{xcolor}

\title{A GIS-Driven Satellite Monitoring System for Industrial Green-Belt Compliance and Change Detection}

\author{\IEEEauthorblockN{Rakshit Verma}
\IEEEauthorblockA{Independent Research and Development\\
Email: rakshit@example.com}
}

\begin{document}
\maketitle

\begin{abstract}
This paper presents a practical monitoring system for pollution-control workflows that integrates industrial-premises boundaries with multi-temporal satellite imagery to detect vegetation decline and potential construction change inside monitored premises. The system combines geometry-aware raster masking, spectral-rule baselines, optional pixel-level trainable models, and automated compliance reporting. A desktop-first user interface was designed to remove direct JSON interaction and support operations teams through human-readable results and guided input validation. Results include before-vs-after metrics, georeferenced violation alerts, and annotated visual evidence panels suitable for regulatory review and periodic monitoring.
\end{abstract}

\begin{IEEEkeywords}
GIS, industrial compliance, remote sensing, Sentinel-2, land-use change, NDVI, environmental monitoring
\end{IEEEkeywords}

\section{Introduction}
Environmental compliance monitoring in industrial areas requires site-specific boundaries, consistent temporal analysis, and auditable evidence generation. Manual inspections are expensive and sparse, while many digital workflows fail due to fragmented boundary data, inconsistent imagery handling, and non-operational reporting outputs. This work addresses these gaps by building an end-to-end monitoring pipeline focused on: (i) premise identification from GIS boundaries, (ii) temporal vegetation and built-up change analysis, (iii) report generation with geo-coordinates and annotated imagery, and (iv) scalable repeat monitoring.

\section{System Objectives}
The implemented system targets five operational objectives:
\begin{enumerate}
    \item Identify factory premises from KGIS/KIADB style boundary inputs.
    \item Analyze multi-date satellite scenes to detect green-cover reduction and construction change candidates.
    \item Apply model-driven vegetation and land-use classification (rule-driven baseline with optional trainable pixel models).
    \item Generate compliance reports with temporal statistics, georeferenced alerts, and visual evidence.
    \item Support continuous, repeatable monitoring runs for board-level operations.
\end{enumerate}

\section{Architecture}
\subsection{Data Ingestion}
Boundary inputs are accepted as zipped shapefiles, shapefiles, and GeoJSON. Premise polygons are rasterized onto the analysis grid using affine transforms and CRS checks.

Satellite scenes are provided as pre-aligned multispectral arrays (blue, green, red, NIR, optional SWIR) in NPZ format. For open-data execution, Sentinel-2 L2A scenes are discovered via STAC and reprojected to EPSG:4326 analysis grids.

\subsection{Monitoring Core}
For each premise and time pair (before, after), the pipeline computes:
\begin{itemize}
    \item Vegetation masks based on NDVI thresholds or optional model inference.
    \item Built-up masks using NDBI-style logic and spectral constraints or optional model inference.
    \item Green-loss and new-construction candidate masks as temporal differencing outputs.
\end{itemize}
Connected regions are extracted and filtered by minimum area and pixel constraints.

\subsection{Model Layer}
Two operation modes are supported:
\begin{itemize}
    \item \textbf{Spectral-rules mode:} deterministic NDVI/NDBI-based classification.
    \item \textbf{ML mode:} trainable pixel logistic models for vegetation and built-up classes, with model bundle save/load for repeat use.
\end{itemize}

\subsection{Reporting Layer}
The reporting module produces:
\begin{itemize}
    \item Compliance JSON and Markdown reports.
    \item Before/after annotated images.
    \item Comparison panels with temporal metadata.
    \item Alert entries containing severity, type, and geo-coordinates.
\end{itemize}

\section{Desktop Operations Interface}
A Tkinter desktop studio was implemented to make the workflow operator-friendly and executable as a packaged desktop application. The interface includes:
\begin{itemize}
    \item Boundary inspection with direct file upload and geometry diagnostics.
    \item OSM candidate discovery with region query and output management.
    \item Compliance analysis form with guided required inputs and no raw JSON display.
    \item Continuous monitoring tab for scheduled repeated runs and history review.
\end{itemize}
Outputs are presented in human-readable summaries and tabular alert views.

\section{Methodology}
\subsection{Premise Masking}
Given affine transform $T$, each polygon is rasterized onto the scene grid. Invalid overlap and CRS mismatch are explicitly rejected to prevent silent geometric errors.

\subsection{Temporal Vegetation Change}
NDVI is computed per pixel:
\begin{equation}
\mathrm{NDVI} = \frac{NIR - Red}{NIR + Red + \epsilon}
\end{equation}
Vegetation masks $V_b$ and $V_a$ for before and after dates produce green-loss mask:
\begin{equation}
L_g = V_b \land \neg V_a
\end{equation}

\subsection{Construction Candidate Detection}
Built-up masks $B_b$ and $B_a$ are estimated from NDBI-style features and thresholds (or learned classifier), then differenced:
\begin{equation}
C_n = B_a \land \neg B_b
\end{equation}
Connected components in $C_n$ are filtered by minimum pixel and area constraints.

\subsection{Compliance Scoring}
Metrics include green-cover percentage before/after, delta in percentage points, green-loss area, and new-construction area. Alerts are triggered when:
\begin{itemize}
    \item Green cover falls below required threshold.
    \item Green-cover reduction exceeds allowed drop.
    \item New construction candidate regions exceed minimum area.
\end{itemize}

\section{Experimental Run Summary}
In the implemented open-data run over Udyambag, Sentinel-2 scenes were selected for two dates and evaluated against an OSM-derived industrial polygon. Generated artifacts included report documents, annotated evidence images, and temporal metrics. The workflow executed end-to-end through the CLI and desktop UI paths.

\section{Scalability and Deployment}
Continuous monitoring is supported through batch iterations and interval-based reruns with persisted history outputs. This supports monthly or quarterly compliance pipelines, trend dashboards, and future scheduler integration.

The desktop-first design supports packaging into a standalone executable for board or field use without requiring API deployment.

\section{Limitations and Future Work}
Current limitations include dependence on NPZ scene packaging and screening-level construction detection at medium resolution. Future improvements include:
\begin{itemize}
    \item Direct GeoTIFF ingestion.
    \item Strongly supervised model training from validated labels.
    \item High-resolution imagery integration for finer construction detection.
    \item Automated report export to board-standard PDF templates.
\end{itemize}

\section{Conclusion}
This work demonstrates a practical and extensible compliance-monitoring system that bridges GIS premise boundaries, satellite time-series analysis, model-enabled classification, and operation-ready reporting. By adding a human-readable desktop interface and repeat-run support, the system is positioned for real-world environmental board workflows.

\begin{thebibliography}{99}
\bibitem{sentinel2}
ESA, ``Sentinel-2 Mission,'' \url{https://www.esa.int/Applications/Observing_the_Earth/Copernicus/Sentinel-2}.

\bibitem{s2sr}
Google Earth Engine, ``COPERNICUS/S2\_SR\_HARMONIZED,'' \url{https://developers.google.com/earth-engine/datasets/catalog/COPERNICUS_S2_SR_HARMONIZED}.

\bibitem{dynamicworld}
Google Earth Engine, ``GOOGLE/DYNAMICWORLD/V1,'' \url{https://developers.google.com/earth-engine/datasets/catalog/GOOGLE_DYNAMICWORLD_V1}.

\bibitem{s1}
ESA, ``Sentinel-1 Mission,'' \url{https://www.esa.int/Applications/Observing_the_Earth/Copernicus/Sentinel-1}.

\bibitem{kiadb}
KIADB GIS Portal, \url{https://kiadb.karnataka.gov.in/kiadbgisportal/}.
\end{thebibliography}

\end{document}
